\section{Discussion}
\subsection{Evidence Delivery and Localization Output}
The joint diagnostics change how aggregate scores should be read. Agent's high annotated-source coverage does not entail coverage of every benchmark role: 80 of its 108 uncovered annotations with source available are EPs, whereas the prompt requests likely vulnerable locations without separately requiring EP and CO. The two counterexamples show why an uncovered entry definition cannot automatically be attributed to failed vulnerability reasoning. Evaluations that require complete entry-to-operation coverage should state that task explicitly.

For RAG, 91 of 96 hits without annotated-source delivery disappear under the width-limited rule. This makes prediction extent central to interpreting apparent localization without source. Titles, descriptions, other code, and model knowledge remain possible information routes; the saved outcomes do not identify their individual contributions. The annotation partition is therefore a joint coverage description, not a causal chain from evidence acquisition to reasoning success.

Schema validity provides another distinct measurement. A syntactically accepted location can name a nonexistent file or exceed its length. The clipped extent makes the returned range interpretable as repository lines while preserving all outcomes, but does not measure actual analyst time. Agent's narrower ranges and higher CO recall describe its returned locations; assessing their practical usefulness also requires checking explanations and task-specific review needs.

\subsection{Completion Is Part of the Measured Configuration}
The full-denominator decomposition quantifies the completion contribution: valid Agent-output pairs contribute positively to the file contrast, but failed pairs contribute a larger negative amount. The observed ordering therefore depends materially on final-output reliability. Keeping both contributions on the original denominator exposes that dependency without selecting a favorable subset as a replacement result.

All 31 Agent failures follow closure of native tool declarations, and 22 valid answers also follow closure. On normal budget exhaustion the implementation disables tools without adding a new model-visible final instruction. Every terminal DSML failure persists after a later JSON correction; empty outputs receive no correction. These observations identify a protocol boundary worthy of intervention, while task difficulty and provider/adapter behavior remain alternative explanations. Estimating performance under a changed completion protocol requires a separately frozen intervention on an appropriately defined sample, including successful exhausted trajectories.

\subsection{Interpreting the Baselines and Budgets}
The frozen Agent--Pack primary contrast measures a query-dependent configuration against a weak lexical prefix. We retain it as registered, while Agent--RAG and the joint coverage diagnostics supply the more informative comparison. Pack describes the limitations of lexical ordering on these snapshots; structural neighborhoods, summaries, or query-driven packing would be different comparators.

Agent's smaller acquisition ledger accompanies much larger cumulative provider input. History replay increases input without delivering new source. The shared ceiling therefore defines one acquisition operating point, with measured recall, output reliability, expense, and time. Matching cumulative input or scanning acquisition limits would address different resource questions. Offline index preparation remains an additional investment for RAG.

\subsection{What the Supplementary Study Establishes}
S1 supplies a shared nonblank candidate universe on the known sample while jointly changing rendering, interaction, schema, ranking details, scheduling, and timing. Its differences from S0 describe those complete configurations. Both studies used dataset-supplied entry titles, so their numbers belong to the title-assisted setting; performance after removing annotation-added title information remains unmeasured.

The empirical contribution is the combination of these distinctions in an inspectable same-model comparison: source-rich outputs can omit an unrequested annotation role, low-source tolerant hits can depend on broad intervals, and schema-valid outputs can contain unusable coordinates. These results explain why source coverage, annotation recall, width, and location validity should be reported together for this task. They characterize the measured configurations without establishing a general ranking of acquisition policies.
